\documentclass[11pt]{article}

\usepackage[margin=1.1in]{geometry}
\usepackage{amsmath}
\usepackage{amssymb}
\usepackage{amsthm}
\usepackage{booktabs}
\usepackage{microtype}
\usepackage[round,authoryear]{natbib}
\usepackage[colorlinks=true,allcolors=blue]{hyperref}

\makeatletter
\renewcommand\paragraph{\@startsection{paragraph}{4}{\z@}%
  {3.25ex \@plus1ex \@minus.2ex}{-1em}%
  {\normalfont\normalsize\itshape}}
\makeatother

\newtheoremstyle{itheads}{\topsep}{\topsep}%
  {\itshape}{}{\itshape}{.}{ }{\thmname{#1}\thmnumber{ #2}%
  \thmnote{ (#3)}}
\theoremstyle{itheads}
\newtheorem{proposition}{Proposition}
\newtheorem{corollary}{Corollary}

\newcommand{\E}{\mathbb{E}}
\newcommand{\PP}{\mathbb{P}}
\DeclareMathOperator{\arcosh}{arcosh}

\title{\textbf{Clay Shooting:\\ Equilibrium Effort Solving Outstanding Problems}}
\author{Peter Cotton}
\date{\today}

\begin{document}
\maketitle

\begin{abstract}
\noindent Two mathematicians at work on the same problem, three years to an
attempt, three hundred thousand dollars a year to keep one of them fed, and
one Millennium problem solved in a hundred and fifty-nine problem-years.
Those four numbers imply that solving one is worth about \$96M to whoever
is attacking it, against a posted purse of \$1M. The inversion runs through
a price list: a Brownian attempt run between a target and an abandonment
floor buys success probability $p$ at expected cost $\zeta^{2}p/(1-p)$,
where $\zeta^{2}$ is the cost of a diffusive traverse of the distance still
to go, and the identity holds at every floor, so an attempt's cost and
success rate pin $\zeta^{2}$ outright. Sizing the attempt against a reward
$A$ gives $p=1-\zeta/\sqrt A$ and $k=p(1-p)A$, and the observed hazard then
fixes $p$, leaving $A=cN/(h(1-p))$ in the cost rate, the number of
attackers and the hazard, with the duration of an attempt cancelling. The
level is not identified by one event: the same arithmetic gives \$18M at the
top of the hazard interval and \$3.8B at the bottom. What is identified is
the \emph{surplus}, since optimal sizing forces
$A-\zeta^{2}=k(2-p)/(1-p)$, about twice the cost of one attempt whatever the
reward, so the purse is $1\%$ of the value and $55\%$ of the surplus, and
that second share moves by one percentage point across the whole interval.
The purse is a rounding error against what a Clay problem is worth and a
majority stake in what makes attacking one worthwhile.
\end{abstract}

\section{What a Millennium problem is worth}

The Clay Mathematics Institute posted one million dollars on each of seven
problems in May 2000 \citep{clay2000}. One has been solved, in 2002--2003,
and the person entitled to the purse declined it. That refusal says the
reward is not denominated in escrow. It does not say what the reward is.

This paper computes it. Suppose two people are at work on a given problem
at any time, an attempt takes about three years, a senior mathematician
costs about three hundred thousand dollars a year all in, and problems fall
at the observed rate of one per hundred and fifty-nine problem-years. Then
solving a Millennium problem is worth about \$96M to the person attacking
it, and the purse is $1.0\%$ of that.

The inversion has one moving part, a price list for persistence. An attempt
run between a target at distance $d$ and an abandonment floor buys success
probability $p$ at expected cost
\begin{equation}
k(p)=\zeta^{2}\,\frac{p}{1-p},
\qquad
\zeta^{2}=\frac{c\,d^{2}}{\sigma^{2}},
\label{eq:intro-price}
\end{equation}
where $\zeta^{2}$ is the cost rate times the diffusive time to cover the
remaining distance. Optimizing $Ap-k(p)$ against a reward $A$ gives
$p=1-\zeta/\sqrt A$ and $k=p(1-p)A$, which inverts to $A=k/(p(1-p))$. The
hazard supplies $p$ through an accounting identity, and what comes out is
\begin{equation}
A=\frac{cN}{h(1-p)},
\qquad
p=\frac{hT}{N},
\end{equation}
in the cost rate $c$, the number $N$ of concurrent attempts, and the hazard
$h$. This is Proposition~\ref{prop:invert}. In English: the duration $T$ of an attempt cancels, so how long people
persist does not enter the level, only how many of them there are and how
fast problems fall.

One event does not identify that level. The exact Poisson interval for the
hazard is a factor of two hundred wide, and it carries $A$ from \$18M to
\$3.8B. Anyone quoting \$96M is quoting the midpoint of an interval that
contains almost anything.

The surplus is a different matter, and it is the paper's point. Optimal
sizing forces
\begin{equation}
A-\zeta^{2}=k\,\frac{2-p}{1-p}\;\ge\;2k,
\label{eq:intro-surplus}
\end{equation}
so the reward exceeds the threshold at which anyone would start by about
twice the cost of a single attempt, however large the reward. Here $T$
survives and $N$ and $h$ drop out. The surplus is \$1.8M, the purse is
$55\%$ of it, and that share moves from $55.5\%$ to $54.1\%$ across the
entire hazard interval. Deleting the purse from the reward costs $55\%$ of
the equilibrium solvers.

So the purse is a rounding error against the value of a Clay problem and a
majority stake in the surplus that makes attacking one worthwhile. The two
ratios differ by a factor of fifty-three because the level and the surplus
scale in disjoint parameters. Dividing a purse by a prize answers the wrong
question.

The rest of the paper earns \eqref{eq:intro-price} to
\eqref{eq:intro-surplus} and then spends them. Section~\ref{sec:price}
proves the price list, whose last clause is the useful one: it holds at
every floor, optimal or not, at any discount rate, so it identifies
$\zeta^{2}$ from an observed attempt without assuming the attempt was well
run. Section~\ref{sec:invert} does the Millennium inversion.
Section~\ref{sec:eq} closes the competitive loop, where the reward is shared
rather than won outright and the equilibrium is a scalar root; in a large
pool it collapses to $G(y^{*})=V_{e}/V$ with $G(y)=(1-e^{-y})/y$, and
competition bids the effective reward down to break-even exactly.
Section~\ref{sec:purse} designs the purse: a funder splitting escrowed
capital posts $1.7564$ times break-even on each problem whatever the
capital, and solves $71.5\%$ of them.

\paragraph{Scope.}
Sections \ref{sec:price}--\ref{sec:purse} treat the lottery allocation:
every entrant pays its attempt to completion and one winner is drawn among
the successes. The first-solution race, in which costs cease when a rival
announces and the rival hazard varies with calendar time, is the companion
note's coupled system \citep{cotton2026races}, and it is not solved here.
The constants below are lottery constants.


\section{The price list}
\label{sec:price}

Fix one researcher and one approach. While working, the researcher
observes $dX_{t}=\sigma\,dW_{t}$ with $\sigma>0$, starting at $x_{0}$.
First passage to a target $b>x_{0}$ is a solution, working costs $c>0$ per
unit time, and abandonment at a floor $a<x_{0}$ is irreversible. Write
$d=b-x_{0}$ for the distance still to travel, $L=b-a$ for the continuation
width, and
\begin{equation}
\zeta^{2}=\frac{c\,d^{2}}{\sigma^{2}}
\end{equation}
for the \emph{hardness}, the cost rate times the diffusive time
$d^{2}/\sigma^{2}$ to cover the remaining distance.

\begin{proposition}[Persistence is bought at a convex price]
\label{prop:price}
For any width $L>d$,
\begin{equation}
p=1-\frac{d}{L},
\qquad
k=\frac{c\,d\,(L-d)}{\sigma^{2}}=\zeta^{2}\,\frac{p}{1-p}.
\label{eq:price}
\end{equation}
The map $k(\cdot)$ is strictly increasing and strictly convex on $[0,1)$,
with $k(0)=0$ and $k(p)\to\infty$ as $p\to1$, and an observed pair
$(k,p)$ determines
\begin{equation}
\zeta^{2}=k\,\frac{1-p}{p}
\label{eq:identify}
\end{equation}
with no reference to the prize, the discount rate, the rival hazard, or
whether $L$ was chosen well.
\end{proposition}

\begin{proof}
Optional stopping applied to $X$ and to $X_{t}^{2}-\sigma^{2}t$ at
$\tau=\tau_{a}\wedge\tau_{b}$, legitimate since $X$ is bounded on $(a,b)$
and $\E[\tau]<\infty$, gives the gambler's-ruin identities
$p=(x_{0}-a)/(b-a)=1-d/L$ and
$\E[\tau]=(x_{0}-a)(b-x_{0})/\sigma^{2}=d(L-d)/\sigma^{2}$, so
$k=c\E[\tau]$. Substituting $L=d/(1-p)$ gives
$k=c\,d^{2}p/\big(\sigma^{2}(1-p)\big)$. Convexity and the limits follow.
\end{proof}

The last clause of Proposition~\ref{prop:price} is what makes the rest
work. Standard race models specify a cost-to-hazard technology and then
compute; here the technology is a theorem about Brownian motion, and it
holds off the optimal path as well as on it.

\begin{corollary}[Optimal attempt size]\label{cor:size}
Let $A>0$ be the expected reward conditional on solving. The attempt
maximizing $Ap-k(p)$ has
\begin{equation}
p=\Big(1-\frac{\zeta}{\sqrt A}\Big)_{+},
\qquad
k=p(1-p)A,
\qquad
\text{value}=Ap^{2}=\big(\sqrt A-\zeta\big)_{+}^{2},
\label{eq:size}
\end{equation}
and is worth mounting exactly when $A>\zeta^{2}$. Expected spend never
exceeds $A/4$, with equality only at $p=1/2$.
\end{corollary}

\begin{proof}
$Ap-\zeta^{2}p/(1-p)$ has derivative $A-\zeta^{2}/(1-p)^{2}$, decreasing
in $p$ and vanishing at $(1-p)^{2}=\zeta^{2}/A$. That stationary point
lies in $(0,1)$ exactly when $A>\zeta^{2}$; otherwise the derivative is
negative throughout and $p=0$. Substituting $1-p=\zeta/\sqrt A$ gives
$k=\zeta p\sqrt A=p(1-p)A$ and $Ap-k=Ap^{2}=(\sqrt A-\zeta)^{2}$. Finally
$k/A=p(1-p)\le1/4$.
\end{proof}

The threshold $\zeta^{2}$ has a reading. A reward must cover an undirected
diffusive crossing of the remaining gap before anyone starts at all, and
the quarter-of-the-reward bound says the expensive attempts are the
genuinely uncertain ones. Cheap long shots and near-certainties both
spend little.

\begin{corollary}[The surplus is a small multiple of one attempt]
\label{cor:surplus}
If the attempt is optimally sized,
\begin{equation}
A-\zeta^{2}=k\,\frac{2-p}{1-p}\;\ge\;2k,
\label{eq:surplus}
\end{equation}
with the bound approached as $p\downarrow0$ and equal to $3k$ at
$p=1/2$.
\end{corollary}

\begin{proof}
$A=\zeta^{2}/(1-p)^{2}$, so
$A-\zeta^{2}=\zeta^{2}p(2-p)/(1-p)^{2}$, and $k=\zeta^{2}p/(1-p)$.
\end{proof}

Corollary~\ref{cor:surplus} is the least expected line in the paper.
However enormous the reward, if attempts are long shots then the reward
barely clears its threshold, and it clears it by roughly twice the cost of
one attempt. The level of the reward and the surplus it generates are
different quantities, and only the second governs behavior.
Section~\ref{sec:invert} is where that distinction earns its keep.

\paragraph{A verification.}
A seeded Euler simulation, given only the boundary location and the
increment law, reproduces \eqref{eq:size}. At $A=100$, $c=\sigma=1$ and
$d=4$, so $\zeta=4$, $L=10$ and $p=0.6$, fifty thousand paths at
$dt=2\times10^{-3}$ give $\hat p=0.6051$ (s.e.\ $0.0022$) and
$\hat k=24.05$ (s.e.\ $0.092$) against $p=0.6$ and $k=24$. The upward bias
in $\hat p$ is Euler overshoot at the upper boundary and shrinks with
$dt$. The price identity \eqref{eq:price} is checked against $k$ computed
directly from the width at five widths, with residuals at $10^{-14}$.

\section{The implied reward}
\label{sec:invert}

Proposition~\ref{prop:price} identifies $\zeta^{2}$ from an attempt's cost
and success rate through \eqref{eq:identify}. Corollary~\ref{cor:size} inverts an optimally sized
attempt to the reward it was chasing. What is missing is $p$, and the
historical record supplies it.

\begin{proposition}[The implied reward]\label{prop:invert}
Let $N$ attempts run per problem at all times, each of mean duration $T$ at
cost rate $c$, each succeeding with probability $p$ and sized optimally
against a reward $A$. If problems are solved at hazard $h$ per
problem-year, then
\begin{equation}
p=\frac{hT}{N},
\qquad
A=\frac{cN}{h(1-p)},
\qquad
\zeta^{2}=\frac{cN(1-p)}{h},
\qquad
A-\zeta^{2}=cT\,\frac{2-p}{1-p}.
\label{eq:invert}
\end{equation}
The level $A$ is free of $T$ up to the factor $(1-p)$; the surplus
$A-\zeta^{2}$ is free of $N$ and $h$, and tends to $2cT$ as $p\downarrow0$.
\end{proposition}

\begin{proof}
Completed attempts arrive at rate $N/T$ per problem-year and each solves
with probability $p$, so $h=pN/T$, which is the first identity. One attempt
costs $k=cT$. Corollary~\ref{cor:size} inverted gives $A=k/(p(1-p))$ and
Proposition~\ref{prop:price} gives $\zeta^{2}=k(1-p)/p$. Substituting
$k=cT$ and $p=hT/N$,
\[
A=\frac{cT}{(hT/N)(1-p)}=\frac{cN}{h(1-p)},
\qquad
\zeta^{2}=\frac{cT(1-p)N}{hT}=\frac{cN(1-p)}{h},
\]
and \eqref{eq:surplus} gives $A-\zeta^{2}=k(2-p)/(1-p)$ with $k=cT$. The limits are immediate.
\end{proof}

The two scalings are disjoint, and that is the whole of what follows. The
level is $cN/h$: what a solution is worth rises with the cost of the people
chasing it and with how many are chasing, and falls with how often problems
fall. The surplus is $2cT$: what makes chasing worthwhile rises with the
cost of one attempt and nothing else. A funder who moves the purse moves
the second and leaves the first alone.

\paragraph{The Millennium numbers.}
One solution since May 2000, with the solved problem exposed only until
2003, gives an exposure of $6\times26+3=159$ problem-years and
$\hat h=0.0063$ per problem-year. At $c=\$300$k per year, $N=2$ and $T=3$
years, \eqref{eq:invert} returns $p=0.0094$, an attempt cost of $\$0.9$M,
$\zeta^{2}=\$94.5$M and
\begin{equation}
A=\$96.3\text{M},
\qquad
A-\zeta^{2}=\$1.81\text{M},
\end{equation}
against $2k=\$1.80$M. The purse is $1.04\%$ of the level and $55.3\%$ of the
surplus, a ratio of fifty-three between the two shares. Dilution is
negligible at these odds: the reward before sharing is $\$96.8$M, since
$\E[(1+K)^{-1}]=0.9953$ with two attackers at $p=0.0094$.

\begin{center}
\begin{tabular}{rrrr}
\toprule
& \multicolumn{3}{c}{implied level $A$, \$M}\\
\cmidrule(l){2-4}
cost rate & $N=1$ & $N=2$ & $N=5$\\
\midrule
\$200k & 32 & 64 & 160\\
\$300k & 49 & 96 & 239\\
\$400k & 65 & 128 & 319\\
\bottomrule
\end{tabular}
\end{center}

\paragraph{The level is not identified and the surplus is.}
The exact Poisson interval for one event in $159$ problem-years is
$[1.6\times10^{-4},\,3.5\times10^{-2}]$, which is to say a solution every
$29$ to $6280$ years. Since $A$ is proportional to $1/h$, that interval is
inherited whole.

\begin{center}
\begin{tabular}{rrrrr}
\toprule
$h$ & $A$, \$M & purse${}/A$ & surplus, \$M & purse${}/$surplus\\
\midrule
$1.6\times10^{-4}$ & 3769 & $0.027\%$ & 1.800 & $55.5\%$\\
$6.3\times10^{-3}$ & 96.3 & $1.04\%$  & 1.809 & $55.3\%$\\
$3.5\times10^{-2}$ & 18.1 & $5.53\%$  & 1.850 & $54.1\%$\\
\bottomrule
\end{tabular}
\end{center}

The level moves by a factor of two hundred and the purse's share of the
surplus moves by one and a half percentage points. This is not a
coincidence to be admired: $h$ enters $A$ and $\zeta^{2}$ the same way and
cancels in the difference, exactly as \eqref{eq:invert} says. The quantity
everyone would quote is the one the record cannot pin down, and the quantity
that governs behavior is the one it can.

\paragraph{What the purse buys.}
Deleting the purse from the reward and re-solving the large-pool equilibrium
of Proposition~\ref{prop:eqlimit} takes expected solvers from $0.0380$ to
$0.0171$, a fall of $55.1\%$, tracking the purse's $55.3\%$ share of the
surplus. Near threshold $y^{*}\approx2(V/V_{e}-1)$, so solvers are
proportional to the excess of the reward over break-even and the purse's
share of that excess is its share of the effort. Across the hazard interval
the fall runs $53\%$ to $56\%$.

Two of the four inputs come from outside the record. The cost rate $c$ and
the number $N$ of concurrent attempts are stipulated, and $A$ is
proportional to both, so the level inherits their error too. The duration
$T$ is stipulated and does not matter for the level. The surplus needs only
$c$ and $T$, and the cost of a serious attempt is the quantity to go and
measure.

Setting $V_{e}=\zeta^{2}$ assumes no cost of entering beyond the flow cost
of working. A setup cost raises $V_{e}$ to $(\zeta+\sqrt K)^{2}$ and shrinks
the excess the purse is measured against, so every share above is a lower
bound on the purse's.

\paragraph{Why researchers stop.}
The remaining Millennium question is not what the reward is but why anyone
abandons it, and here the arithmetic belongs to the companion note's
first-solution race rather than the lottery, since it needs the option
destroyed rather than diluted. There an attempt faces $\rho=r+h$ and
competition supplies a share $h/(r+h)$ of the impatience. Two adjustments
push the same way: the problem-level hazard is the aggregate, including the
focal researcher, whose own contribution must come out, and with $N$
symmetric attackers the rival hazard is $(N-1)/N$ of it; and any positive
discount rate enters the denominator.

\begin{center}
\begin{tabular}{rccc}
\toprule
$r$ & aggregate share & rival share, $N=3$ & aggregate at interval upper\\
\midrule
0.03 & 0.173 & 0.123 & 0.539\\
0.05 & 0.112 & 0.077 & 0.412\\
0.10 & 0.059 & 0.040 & 0.260\\
\bottomrule
\end{tabular}
\end{center}

A solution hazard of order $10^{-2}$ per year is small beside any plausible
discount rate, so rivalry is unlikely to be the main reason anyone abandons
a Millennium problem. They stop because the shells cost money and the years
pass.

\paragraph{When the metaphor of the title applies.}
All the interesting variation sits in $\zeta=d\sqrt c/\sigma$. An attempt
run at machine cost lowers $c$, and a genuinely broader search raises
$\sigma$; both lower $\zeta$, which lowers the break-even purse
quadratically and raises expected solvers. If that compresses expected
time-to-solution from decades to a few years, $h$ rises by two orders of
magnitude, and by \eqref{eq:invert} the implied level falls in proportion
while the rivalry share moves from under a tenth to above four fifths. Only
then is the purse a clay pigeon, thrown up for several guns with one hit
counting. This is a graded, falsifiable prediction.


\section{Equilibrium attempt size}
\label{sec:eq}

Now the contest. There are $n$ researchers, a prize $V$, and a common
hardness $\zeta^{2}$. Each researcher chooses a width, equivalently a
success probability $p_{i}\in[0,1)$, and pays $k(p_{i})$; choosing
$p_{i}=0$ is staying out. Successes are independent, and if several
researchers succeed one is drawn uniformly for the prize.

Conditional on solving, researcher $i$'s expected reward is
\begin{equation}
A_{i}=V\,\E\Big[\frac{1}{1+K_{-i}}\Big]
=V\int_{0}^{1}\prod_{j\ne i}\big(1-p_{j}+p_{j}z\big)\,dz,
\label{eq:Ai}
\end{equation}
by $1/(m+1)=\int_{0}^{1}z^{m}dz$, where $K_{-i}$ counts successful rivals.
Payoffs are $A_{i}p_{i}-k(p_{i})$, so Corollary~\ref{cor:size} makes each
best response $p_{i}=(1-\zeta/\sqrt{A_{i}})_{+}$. Rivals enter through one
number. In the symmetric case \eqref{eq:Ai} collapses to
\begin{equation}
A(p)=V\,\frac{1-(1-p)^{n}}{np}.
\label{eq:Asym}
\end{equation}

\begin{proposition}[The symmetric equilibrium is a scalar root]
\label{prop:eqfinite}
$A$ in \eqref{eq:Asym} is strictly decreasing in $p$ on $(0,1]$, with
$A(0^{+})=V$ and $A(1)=V/n$, and strictly decreasing in $n$. If
$V\le\zeta^{2}$ the only symmetric equilibrium is $p=0$. If
$V>\zeta^{2}$ there is a unique symmetric equilibrium $p_{n}\in(0,1)$, the
root of
\begin{equation}
p=1-\frac{\zeta}{\sqrt{A(p)}},
\end{equation}
it decreases in $n$, and each researcher earns $A(p_{n})p_{n}^{2}>0$.
\end{proposition}

\begin{proof}
Write $w=1-p$. Then $\partial_{p}\big[(1-w^{n})/p\big]$ has the sign of
$g(p)=npw^{n-1}-1+w^{n}$, and $g(0)=0$ with
$g'(p)=-n(n-1)pw^{n-2}<0$, so $g<0$ and $A$ decreases in $p$. For
monotonicity in $n$, the difference
$(n+1)(1-w^{n})-n(1-w^{n+1})=1-(n+1)w^{n}+nw^{n+1}$ vanishes at $w=1$ with
derivative $n(n+1)w^{n-1}(w-1)<0$, hence is positive for $w<1$, giving
$A_{n}>A_{n+1}$.

Set $\Phi(p)=(1-\zeta/\sqrt{A(p)})_{+}-p$, continuous and strictly
decreasing on $[0,1]$: where the positive part is active it is
$1-\zeta/\sqrt{A(p)}-p$ with $A$ strictly decreasing, and where it is
clamped it is $-p$. Now $\Phi(0)=1-\zeta/\sqrt V$ and
$\Phi(1)=-\zeta\sqrt n/\sqrt V<0$. If $V\le\zeta^{2}$ then $\Phi(0)\le0$
and $p=0$ is the unique root, and it is an equilibrium, since a deviator
faces $A=V\le\zeta^{2}$ and has best response $0$. If $V>\zeta^{2}$ then
$\Phi(0)>0>\Phi(1)$ and there is exactly one root, decreasing in $n$
because $A_{n}$ is.
\end{proof}

\begin{proposition}[Large pool]\label{prop:eqlimit}
Let $n\to\infty$ with $np\to y$. Then $A\to V\,G(y)$ with
$G(y)=(1-e^{-y})/y$, strictly decreasing from $1$ to $0$, and the
equilibrium condition becomes
\begin{equation}
\boxed{\;G(y^{*})=\frac{\zeta^{2}}{V}\;}
\label{eq:limit}
\end{equation}
with a unique positive root whenever $V>\zeta^{2}$. At that root the
expected number of solvers is $y^{*}$, the solve probability is
$1-e^{-y^{*}}$, the effective reward collapses to $A^{*}=\zeta^{2}$,
individual effort and payoff vanish, and aggregate effort equals the
expected payout,
\begin{equation}
\lim n\,k(p_{n})=y^{*}\zeta^{2}=V\big(1-e^{-y^{*}}\big).
\end{equation}
With a setup cost $K>0$ payable on entry,
\begin{equation}
G(y^{*})=\frac{V_{e}}{V},
\quad
V_{e}=\big(\zeta+\sqrt K\big)^{2},
\quad
p^{*}=\frac{\sqrt K}{\zeta+\sqrt K},
\quad
N^{*}=y^{*}\,\frac{\zeta+\sqrt K}{\sqrt K},
\end{equation}
where $V_{e}$ is the purse at which a lone researcher exactly breaks even,
$A^{*}=V_{e}$, each attempt costs $\zeta\sqrt K$, and $N^{*}\to\infty$ as
$K\downarrow0$.
\end{proposition}

\begin{proof}
$(1-y/n)^{n}\to e^{-y}$ gives $A\to VG(y)$. Since $p=1-\zeta/\sqrt A$ and
$p\to0$, we need $\zeta/\sqrt{A^{*}}\to1$, so $A^{*}=\zeta^{2}$ and
\eqref{eq:limit} follows; uniqueness is strict monotonicity of $G$. Then
$k=p(1-p)A\to0$ and the payoff $Ap^{2}\to0$, while
$nk\to y\,A^{*}=y^{*}\zeta^{2}=V\,y^{*}G(y^{*})=V(1-e^{-y^{*}})$. With a
setup cost, free entry requires $Ap^{2}=K$, and substituting
$p=1-\zeta/\sqrt A$ gives $(\sqrt A-\zeta)^{2}=K$, so
$\sqrt{A^{*}}=\zeta+\sqrt K$. Then $p^{*}=1-\zeta/(\zeta+\sqrt K)$,
$k=p^{*}(1-p^{*})A^{*}=\zeta\sqrt K$ and $N^{*}=y^{*}/p^{*}$. Finally
$V_{e}$ solves $(\sqrt V-\zeta)^{2}=K$, the lone researcher's break-even
condition in Corollary~\ref{cor:size}.
\end{proof}

\begin{center}
\begin{tabular}{rrrr}
\toprule
$V/V_{e}$ & $y^{*}$ & $\PP(\text{solved})$ & $n p_{n}$ at $n=10^{4}$\\
\midrule
1.02 & 0.03974 & 0.0390 & 0.03972\\
1.5  & 0.87422 & 0.5828 & 0.87387\\
2    & 1.59362 & 0.7968 & 1.59286\\
4    & 3.92069 & 0.9802 & 3.91742\\
10   & 9.99955 & 0.99995& 9.97959\\
\bottomrule
\end{tabular}
\end{center}

The last column solves Proposition~\ref{prop:eqfinite} exactly at
$n=10^{4}$. A deviation certificate backs the fixed points: at three
computed equilibria, no point of a two-million-point grid over own success
probability beats the equilibrium payoff against equilibrium rivals, the
best falling short by $4\times10^{-12}$ or less. The grid does not know the
closed-form best response, so the check can fail for the right reason.

\paragraph{Solutions are pinned, attempts are not.}
Both $y^{*}$ and the solve probability depend only on $V/V_{e}$. The
entrant count $N^{*}=y^{*}(\zeta+\sqrt K)/\sqrt K$ does not, and it
diverges as the setup cost vanishes, since a great many researchers each
buy a vanishing probability. A claim about how many people attempt a
famous problem is a claim about setup costs.

\paragraph{Competition is competed away exactly.}
$A^{*}=V_{e}$. In a large pool the effective reward equals the purse at
which a lone researcher would exactly break even, so every researcher
behaves as though facing no competition and a prize of $V_{e}$, and
aggregate effort equals the expected payout. This is the classical
free-entry dissipation result with the cost curve derived rather than
posited.

\section{The solo width is not a best response}
\label{sec:committed}

A tempting shortcut fixes attempt lengths at the solo-optimal
$L=\sigma\sqrt{V/c}$ and lets researchers choose only whether to enter.
The shortcut is false, and the failure is exact.

Take $V=100$, $c=\sigma=1$, $d=5$, and two entrants at the solo width
$L=10$, so $p=1/2$ each. Then $A=V(1-p/2)=75$, and the payoff at the solo
width is $75\cdot\tfrac12-25=12.5$. The best response is
$L=\sigma\sqrt{A/c}=\sqrt{75}\approx8.6603$, worth $13.3975$, a gain of
$7.2\%$. Competition shortens attempts, because the prize is diluted and a
diluted prize does not justify the same persistence.

The shortcut also misreports how many people attempt. Holding
$(p,k=p(1-p)V)$ exogenous, any three designated entrants form a pure Nash
equilibrium whenever $p<2-\sqrt2\approx0.5858$: each incumbent earns
$V\big[1-(1-p)^{3}\big]/3-p(1-p)V=Vp^{3}/3>0$, while a fourth would earn
$Vp^{2}\big(-\tfrac12+p-\tfrac{p^{2}}{4}\big)<0$. Two entrants is never a
pure equilibrium, since a third always gains $Vp^{3}/3$ by joining.

The companion note's exact potential selects three entrants on exactly the
same range, and not by coincidence. An exact potential moves by the
deviator's own payoff, so its maximiser leaves three at the $p$ where the
fourth entrant breaks even, which is again $2-\sqrt2$. The committed
increments agree with the entrant payoffs to $2\times10^{-16}$.

The symmetric mixed equilibrium of the same game gives
$nq\to2n/(n-1)$ as $p\downarrow0$, which is $2.667$ at $n=4$ and $2.041$
at $n=50$, reaching two only in the pool limit. Any headline count from
the shortcut is an equilibrium selection in disguise.

\section{Purse design}
\label{sec:purse}

A funder commits escrowed capital $B$ to identical problems with
break-even purse $V_{e}$, choosing the number $M$ of problems and posting
$V=B/M$ on each. Free entry then sets each problem's solve probability by
Proposition~\ref{prop:eqlimit}. Write $H(y)=1-e^{-y}$, so $G=H/y$.

\begin{proposition}[The optimal purse is a fixed multiple of break-even]
\label{prop:purse}
Under Proposition~\ref{prop:eqlimit}, $V=V_{e}\,y/H(y)$, so
$M=B\,H(y)/(V_{e}y)$ and the expected number of problems solved is
\begin{equation}
\Psi=M\,H(y)=\frac{B}{V_{e}}\cdot\frac{H(y)^{2}}{y}.
\end{equation}
Over the continuous relaxation of $M$, $\Psi$ has a unique positive
maximizer, characterized by
\begin{equation}
\boxed{\;e^{y}=1+2y\;}
\end{equation}
at which, exactly,
\begin{equation}
\frac{V^{*}}{V_{e}}=y^{*}+\frac12,
\qquad
\PP(\text{solved})=\frac{2y^{*}}{1+2y^{*}},
\qquad
\Psi^{*}=\frac{B}{V_{e}}\cdot\frac{H(y^{*})^{2}}{y^{*}}.
\end{equation}
Numerically $y^{*}=1.25643121$, $V^{*}=1.75643121\,V_{e}$,
$\PP(\text{solved})=0.71533186$, and
$\Psi^{*}=0.40726438\,B/V_{e}$ on $M^{*}=0.56933627\,B/V_{e}$ problems.
\end{proposition}

\begin{proof}
$G(y)=V_{e}/V$ gives $V=V_{e}y/H(y)$, hence $M$ and $\Psi$. Maximizing
$\log\Psi=2\log H-\log y$ up to a constant gives $2H'/H=1/y$, and with
$H'=e^{-y}$ this is $2ye^{-y}=1-e^{-y}$, that is $e^{-y}(2y+1)=1$, that is
$e^{y}=1+2y$. The function $e^{y}-1-2y$ is convex with value $0$ and
derivative $-1$ at the origin and tends to infinity, so it has exactly one
positive zero. At that zero $H=1-1/(1+2y)=2y/(1+2y)$, so
$y/H=(1+2y)/2=y+\tfrac12$.
\end{proof}

The purse multiple is independent of $B$: capital chooses how many
problems to post, not how much to post on each. Running the same program
through the committed-width shortcut of Section~\ref{sec:committed}, where
$V=V_{\min}(y/H)^{2}$, gives $e^{y}=1+\tfrac32y$, $y^{*}=0.76268856$,
$V^{*}=\big(y^{*}+\tfrac23\big)^{2}V_{\min}=2.04305637\,V_{\min}$ and
$\PP(\text{solved})=0.53358923$. Both constants agree with a
two-million-point brute-force maximization of the exact objective. For a
purse exponent $q$, meaning $V=V_{e}(y/H)^{q}$, the condition is
$e^{y}=1+(1+1/q)y$ and the optimum is
$V^{*}/V_{e}=\big(y^{*}+q/(q+1)\big)^{q}$. Widths that adjust give $q=1$ and
widths that are committed give $q=2$.

Integer $M$ reintroduces budget-dependent rounding, so budget independence
belongs to the relaxation. And expected count solved is one objective
among several: a funder who cares most about the single hardest problem
should not spread at all.

\paragraph{Near threshold the purse is powerful.}
For an attempt sized against $V$, $p=1-\zeta/\sqrt V$, so
$dp/d\log V=(1-p)/2$ and doubling $V$ sends $p\mapsto1-(1-p)/\sqrt2$. From
$p=0.01$ that is $p=0.29996$, and a tenfold purse gives $0.68694$.
Sensitivity is largest where $p$ is smallest, because a small $p$ means
$V$ only just exceeds $\zeta^{2}$. Equation~\eqref{eq:limit} inherits the
steepness, which is why the purse carried $55\%$ of the Millennium effort
in Section~\ref{sec:invert} on a $1\%$ share of the level, and why
Proposition~\ref{prop:purse} places the optimum at $1.76$ times break-even,
inside the steep region by construction.

Far above threshold the picture reverses. The companion note's
continuation width at discount-plus-hazard $\rho>0$ is
$L=\sigma(2\rho)^{-1/2}\arcosh(1+\rho V/c)$, which grows like $\log V$
once $\rho V\gg c$, against $\sigma\sqrt{V/c}$ at $\rho=0$. Purses buy
persistence logarithmically out there, and breadth beats depth. Neither
regime dominates, and $V/V_{e}$ selects which one a contest is in.

\section{Assumptions}

\paragraph{The lottery and the race.}
Sections \ref{sec:price}--\ref{sec:purse} solve the lottery. In the
first-solution race costs cease when a rival announces, the rival hazard
is a function of calendar time, and the equilibrium is the companion
note's coupled stopping and forward-distribution system. The race destroys
the option where the lottery only dilutes it, so the constants here do not
transfer.

\paragraph{Asymmetric equilibria.}
Proposition~\ref{prop:eqfinite} gives uniqueness inside the symmetric
class, and it needs only that $A$ decreases. Outside that class the
question is open. The tempting route is Banach: the symmetric
best-response map has modulus
$|T'|=\zeta|A'(p)|/(2A^{3/2})$, and a modulus below one would settle
uniqueness outright. It is $0.18$ at $n=2$ and $0.77$ at $n=10$, then
$12.0$ at $n=100$ and $500.3$ at $n=10^{4}$, so the argument fails exactly
in the large pool the paper spends most of its time in. Best responses do
fall in rivals' success probabilities, which makes this a game of strategic
substitutes, and that is all we use.

\paragraph{Pausing.}
The companion note's intensity extension permits $u=0$, and one might
expect a pause to soften the walk-away threshold. It does not. With
progress frozen while idle, constant discounting and rival hazard, and no
information arriving, the value function decays at rate $\rho$ while
nothing is gained, so waiting strictly destroys value and the optimal
control is bang-bang. A mathematician who returns to a problem after
twenty years needs something to have changed while it sat: new tools
lowering $c$, or a broader method raising $\sigma$. That is a model in
which $\zeta$ drifts down, which is the mechanism of
Section~\ref{sec:invert}.

\paragraph{Independence.}
The product survival formula behind \eqref{eq:Ai} and behind the rivalry
share needs independence across researchers. In mathematics the dominant
uncertainty is shared: whether the statement is true, whether the method
can work at all. A rival's visible failure is evidence about one's own
prospects, and pricing that requires filtering and a correlated state.

\paragraph{Public progress and partial credit.}
Preprints and talks reveal rival states, which makes a joint-state game.
Partial results are published and rewarded, so an all-or-nothing
first-passage payoff overstates the waste. Verification takes years, the
Poincar\'e purse being paid seven after the preprints, so a claimed proof
that later fails restarts the race and the hazard conflates a claim with a
solution.

\paragraph{Measurement.}
Of the four inputs to Proposition~\ref{prop:invert} only the hazard comes
from data, and it comes from one event. The cost rate and the number of
concurrent attackers are stipulated, and the implied level is proportional
to both. Measuring what a serious attempt costs is harder than any
mathematics here.

\section{Related work}

\paragraph{Racing for a prize.}
\citet{loury1979} and \citet{leewilde1980} set the pattern: a flow of R\&D
spending buys a Poisson hazard of success, and free entry dissipates the
rents. \citet{dasguptastiglitz1980} added industrial structure,
\citet{reinganum1981} the dynamic game, \citet{fudenberg1983} preemption
and leapfrogging. \citet{harrisvickers1985,harrisvickers1987} put the race
on a state space with observable positions and asked who gives up. In all
of them the cost-to-hazard technology is specified;
Proposition~\ref{prop:price} derives it, which is what turns
Section~\ref{sec:eq} into a scalar fixed point.

\paragraph{Private progress and preemption.}
\citet{hopenhaynsquintani2011} is the closest neighbor, with privately
observed progress, an option to stop, and preemption. \citet{bbm2017}
model publishing early against polishing under priority risk.
\citet{moscarinisquintani2010} and \citet{krc2005} study experimentation
when a rival's information is itself a motive to continue. Those papers
close the loop in the harder allocation and leave the cost curve in
reduced form; this one closes the curve in the easier allocation.

\paragraph{Contest design.}
\citet{tullock1980} and \citet{baye1996} price rent dissipation.
\citet{taylor1995} and \citet{chegale2003} design research tournaments,
\citet{fullerton1999} auctions entry to select contestants, and
\citet{halac2017} design contests when the technology must be discovered.
Proposition~\ref{prop:purse} uses a cruder instrument than any of these,
namely how many purses and how large, and returns a fixed multiple of
break-even.

\paragraph{Entry games and potentials.}
\citet{morgan2012} and \citet{fujiaolu2015} study endogenous entry with
subsequent effort. The companion note's fixed-attempt lottery and its
exact potential sit in the tradition of \citet{rosenthal1973} and
\citet{monderershapley1996}, and Section~\ref{sec:committed} uses that
potential to locate the three-entrant equilibria.

\paragraph{Budgeted races.}
The forward-backward architecture of the companion note comes from the
budgeted Brownian race \citep{cotton2026wavefront}, where one controller
prunes many paths under a path-time budget and is paid the terminal
maximum. There a shadow price closes the loop. Here it is a consistency
condition between stopping rules and the reward they share, so the
concavity certificates of that problem say nothing about equilibrium in
this one.

\section{Closing}

An attempt on a hard problem has a price list, $k(p)=\zeta^{2}p/(1-p)$,
and the price list carries most of the paper. It makes the cost curve a
theorem rather than a specification, reduces the competitive equilibrium
to a scalar root in which the effective reward is competed down exactly to
break-even, and delivers an optimal purse at $1.7564$ times that level for
any capital.

Neither breadth nor depth dominates. Below about twice break-even the
marginal dollar buys success probability at rate $(1-p)/2$ per log dollar
and the purse is the binding instrument; far above it, persistence grows
like $\log V$ and the same dollar buys another problem instead. The
boundary is drawn at $V/V_{e}$, and the optimal design sits at $1.76$,
just inside the steep side.

The same dial appears elsewhere. A controller allocating test-time compute
across parallel rollouts faces $\zeta$ as the cost of a traverse and the
identical convex price for reliability, which is why the budgeted race
\citep{cotton2026wavefront} and this contest share an architecture and
differ only in what closes the loop. We leave to future work the existence
and selection theory for the first-solution race, whose allocation rule is
the one the title describes and which nobody has solved.

\begin{thebibliography}{99}

\bibitem[Baye et al.(1996)]{baye1996}
Baye, M. R., Kovenock, D., and de Vries, C. G. (1996).
The all-pay auction with complete information.
\emph{Economic Theory}, 8(2):291--305.

\bibitem[Bobtcheff et al.(2017)]{bbm2017}
Bobtcheff, C., Bolte, J., and Mariotti, T. (2017).
Researcher's dilemma.
\emph{Review of Economic Studies}, 84(3):969--1014.

\bibitem[Che and Gale(2003)]{chegale2003}
Che, Y.-K., and Gale, I. (2003).
Optimal design of research contests.
\emph{American Economic Review}, 93(3):646--671.

\bibitem[Clay Mathematics Institute(2000)]{clay2000}
Clay Mathematics Institute (2000).
The Millennium Prize Problems.
Rules and problem descriptions, \texttt{claymath.org}.
Accessed September 2026.

\bibitem[Cotton(2026a)]{cotton2026wavefront}
Cotton, P. (2026a).
Wavefront pruning in budgeted Brownian races.
Working draft, \texttt{brownianbandit.microprediction.org}.

\bibitem[Cotton(2026b)]{cotton2026races}
Cotton, P. (2026b).
Research prizes with uncertain compute costs: from entry games to
Brownian stopping races.
Working note, \texttt{brownianbandit.microprediction.org}.

\bibitem[Dasgupta and Stiglitz(1980)]{dasguptastiglitz1980}
Dasgupta, P., and Stiglitz, J. (1980).
Uncertainty, industrial structure, and the speed of R\&D.
\emph{Bell Journal of Economics}, 11(1):1--28.

\bibitem[Fu et al.(2015)]{fujiaolu2015}
Fu, Q., Jiao, Q., and Lu, J. (2015).
Contests with endogenous entry.
\emph{International Journal of Game Theory}, 44(2):387--424.

\bibitem[Fudenberg et al.(1983)]{fudenberg1983}
Fudenberg, D., Gilbert, R., Stiglitz, J., and Tirole, J. (1983).
Preemption, leapfrogging and competition in patent races.
\emph{European Economic Review}, 22(1):3--31.

\bibitem[Fullerton and McAfee(1999)]{fullerton1999}
Fullerton, R. L., and McAfee, R. P. (1999).
Auctioning entry into tournaments.
\emph{Journal of Political Economy}, 107(3):573--605.

\bibitem[Halac et al.(2017)]{halac2017}
Halac, M., Kartik, N., and Liu, Q. (2017).
Contests for experimentation.
\emph{Journal of Political Economy}, 125(5):1523--1569.

\bibitem[Harris and Vickers(1985)]{harrisvickers1985}
Harris, C., and Vickers, J. (1985).
Perfect equilibrium in a model of a race.
\emph{Review of Economic Studies}, 52(2):193--209.

\bibitem[Harris and Vickers(1987)]{harrisvickers1987}
Harris, C., and Vickers, J. (1987).
Racing with uncertainty.
\emph{Review of Economic Studies}, 54(1):1--21.

\bibitem[Hopenhayn and Squintani(2011)]{hopenhaynsquintani2011}
Hopenhayn, H., and Squintani, F. (2011).
Preemption games with private information.
\emph{Review of Economic Studies}, 78(2):667--692.

\bibitem[Keller et al.(2005)]{krc2005}
Keller, G., Rady, S., and Cripps, M. (2005).
Strategic experimentation with exponential bandits.
\emph{Econometrica}, 73(1):39--68.

\bibitem[Lee and Wilde(1980)]{leewilde1980}
Lee, T., and Wilde, L. L. (1980).
Market structure and innovation: a reformulation.
\emph{Quarterly Journal of Economics}, 94(2):429--436.

\bibitem[Loury(1979)]{loury1979}
Loury, G. C. (1979).
Market structure and innovation.
\emph{Quarterly Journal of Economics}, 93(3):395--410.

\bibitem[Monderer and Shapley(1996)]{monderershapley1996}
Monderer, D., and Shapley, L. S. (1996).
Potential games.
\emph{Games and Economic Behavior}, 14(1):124--143.

\bibitem[Morgan et al.(2012)]{morgan2012}
Morgan, J., Orzen, H., and Sefton, M. (2012).
Endogenous entry in contests.
\emph{Economic Theory}, 51:435--463.

\bibitem[Moscarini and Squintani(2010)]{moscarinisquintani2010}
Moscarini, G., and Squintani, F. (2010).
Competition and experimentation with private information.
\emph{Journal of Economic Theory}, 145(2):639--660.

\bibitem[Reinganum(1981)]{reinganum1981}
Reinganum, J. F. (1981).
Dynamic games of innovation.
\emph{Journal of Economic Theory}, 25(1):21--41.

\bibitem[Rosenthal(1973)]{rosenthal1973}
Rosenthal, R. W. (1973).
A class of games possessing pure-strategy Nash equilibria.
\emph{International Journal of Game Theory}, 2:65--67.

\bibitem[Taylor(1995)]{taylor1995}
Taylor, C. R. (1995).
Digging for golden carrots: an analysis of research tournaments.
\emph{American Economic Review}, 85(4):872--890.

\bibitem[Tullock(1980)]{tullock1980}
Tullock, G. (1980).
Efficient rent seeking.
In J. M. Buchanan, R. D. Tollison, and G. Tullock, eds.,
\emph{Toward a Theory of the Rent-Seeking Society}.
Texas A\&M University Press.

\end{thebibliography}

\end{document}
